\pdfoutput=1
\documentclass[11pt]{amsart}

\usepackage[margin=1.1in]{geometry}
\usepackage{amsmath,amssymb,amsthm,mathtools}
\usepackage{array,booktabs}
\usepackage{enumitem}
\usepackage{microtype}
\usepackage{hyperref}
\usepackage[nameinlink,capitalize]{cleveref}

\hypersetup{colorlinks=true,linkcolor=blue,citecolor=blue,urlcolor=blue}
\raggedbottom
\numberwithin{equation}{section}

\theoremstyle{plain}
\newtheorem{theorem}{Theorem}[section]
\newtheorem{proposition}[theorem]{Proposition}
\newtheorem{lemma}[theorem]{Lemma}
\newtheorem{corollary}[theorem]{Corollary}

\theoremstyle{definition}
\newtheorem{definition}[theorem]{Definition}

\theoremstyle{remark}
\newtheorem{remark}[theorem]{Remark}

\newcommand{\R}{\mathbb R}
\newcommand{\Wloc}{\mathcal W}
\newcommand{\Diss}{\mathcal D}
\newcommand{\Err}{\mathcal E}
\newcommand{\EntropyGrad}{\mathfrak G}
\newcommand{\State}{\mathfrak S}
\newcommand{\SDiff}{\mathrm{SDiff}}
\newcommand{\diver}{\operatorname{div}}
\newcommand{\tr}{\operatorname{tr}}
\newcommand{\supp}{\operatorname{supp}}
\newcommand{\esssup}{\operatorname*{ess\,sup}}

\title[Localized entropy and compactification]
{A Localized Entropy Functional and State-Space Compactification \\
for Navier--Stokes Singular Branches}

\author{Guillem Duran Ballester}
\email{guillem@fragile.tech}
\date{}

\subjclass[2020]{35Q30, 35B44, 35B40, 76D05}
\keywords{Navier--Stokes equations, singularity formation, local energy estimates, entropy methods, ancient solutions, state-space compactification}

\begin{document}

\begin{abstract}
We introduce a localized entropy formalism for the singular-branch state space
arising in the companion Navier--Stokes manuscripts on local Type~I reduction,
residual Type~I exclusion, and Type~II regularity.  The functional is defined
on bounded centered ancient profiles and on repaired-gauge Type~II branches by
combining local enstrophy, mean-free local pressure, scale-critical velocity
mass, and a normalized backward adjoint weight.  Its derivative is local: after
differentiation, the only error terms are the harmonic-pressure, cutoff, and
rerouting defects already controlled in the companion arguments.  Consequently,
every hypothetical singular branch admits an entropy-compactifying subsequence.
We show that every entropy-critical limit belongs, modulo the rerouting
theorems of the companion papers, to one of five terminal model classes:
compact single-core, multiple-core, rotational or stationary-hull,
scale-cascade, and diffuse or tail-defect.  The exclusion theorems proved in
the companion papers therefore yield a finite entropy-based reduction of the
singularity problem.  The purpose of the present paper is not to replace the
local estimates, but to organize them through a single monotone quantity stated
in standard PDE language.
\end{abstract}

\maketitle
\tableofcontents

\section{Introduction}

The local study of possible finite-time singularities for the three-dimensional
incompressible Navier--Stokes equations
\begin{equation}\label{eq:NS}
	\partial_t u - \Delta u + (u\cdot\nabla)u + \nabla p = 0,
	\qquad \diver u = 0,
\end{equation}
is governed by a persistent tension between scale-critical compactness and the
nonlocal structure of the pressure.  Since Leray's construction of global weak
solutions \cite{Leray1934}, the basic local framework has been provided by the
theory of suitable weak solutions and the partial regularity results of
Scheffer, Caffarelli--Kohn--Nirenberg, Lin, and subsequent authors
\cite{Scheffer1976,CKN1982,Lin1998,GustafsonKangTsai2007}.  In that framework,
the relevant local quantities are the scale-invariant velocity and pressure
functionals, the local energy inequality, and the compactness mechanisms that
survive parabolic rescaling.  The pressure term is indispensable but nonlocal,
and this is precisely what complicates a clean compactification of singular
branches.

The three companion manuscripts
\cite{BallesterSetup,BallesterResidual,BallesterTypeII} resolve this difficulty
by a detailed local state-space stratification.  The setup paper
\cite{BallesterSetup} reduces a singular branch either to a bounded centered
ancient profile obtained by Seregin-type extraction or to a local Type~II
branch.  The residual paper \cite{BallesterResidual} closes the remaining
Type~I residual states.  The Type~II paper \cite{BallesterTypeII} introduces a
repaired-gauge representation, pressure decomposition, Caccioppoli estimates,
multibubble and scale-collapse alternatives, carrier routing, and terminal
profile analysis.  Taken together, these papers yield a finite decision tree.

The purpose of the present paper is to reformulate that tree through a single
localized entropy quantity.  The motivation comes from two classical sources.
First, following Arnold's formal interpretation of Euler flow as geodesic
motion on the group of volume-preserving diffeomorphisms
\cite{Arnold1966,EbinMarsden1970}, one may regard the convective term as a
transport term built into the ambient geometry.  Second, following Hamilton and
Perelman \cite{Hamilton1982,Perelman2002}, one expects a useful organizing
quantity to couple a backward adjoint weight to the dissipative part of the
equation and to penalize failure of compactification.  In the present setting,
however, the functional must be strictly local and compatible with the actual
PDE estimates available in the companion papers.  For that reason, we work with
local enstrophy, local scale-critical velocity mass, and mean-free local
pressure, rather than with a global geometric invariant.

The resulting functional is not presented as an independent replacement for the
companion arguments.  It is an entropy-level restatement of the same local
mechanism.  The monotonicity identity records exactly the defects that the
companion papers already control: local pressure remainders, cutoff errors,
noncompact carrier terms, and the branch changes that arise in the local
state-space.  Once these are accounted for, the singularity analysis reduces to
five terminal model classes.

The main conclusion of the paper is the following.

\begin{theorem}[Entropy reduction of the companion program]
\label{thm:intro-main}
Within the singular-branch framework established in
\cite{BallesterSetup,BallesterResidual,BallesterTypeII}, there exists a
localized entropy functional \(\Wloc\) with the following properties.
\begin{enumerate}[label=\textup{(\roman*)},leftmargin=2.2em]
	\item On every admissible branch state produced by the companion papers,
	\(\Wloc\) satisfies a monotonicity identity up to an error term controlled by
	the local pressure reconstruction, compactness, and routing estimates already
	proved in those papers.
	\item After the local error terms have been absorbed, the vanishing of the
	entropy gradient yields a weighted stationary system; its solutions are,
	modulo the companion rerouting theorems, exactly the five terminal model
	classes used in the state-space compactification.
	\item Every hypothetical singular branch has an entropy-critical
	subsequential limit.
	\item Every entropy-critical limit belongs, modulo the rerouting theorems of
	the companion papers, to one of five terminal model classes:
	compact single-core, multiple-core, rotational or stationary-hull,
	scale-cascade, and diffuse or tail-defect.
	\item The exclusion theorems proved in
	\cite{BallesterSetup,BallesterResidual,BallesterTypeII} eliminate each of
	these five terminal classes.
\end{enumerate}
\end{theorem}

\begin{remark}
Theorem~\ref{thm:intro-main} is an organizational theorem.  The local PDE
content remains in the companion papers.  What is new here is the entropy
formulation: all branch changes are encoded as terms in a local monotonicity
identity, and all terminal mechanisms are grouped into five model classes.
\end{remark}

\subsection*{Relation to the literature}

The local Navier--Stokes background used throughout consists of Leray's weak
solutions \cite{Leray1934}, the suitable weak solution framework and partial
regularity results \cite{Scheffer1976,CKN1982,Lin1998}, classical critical-space
and local theory \cite{Ladyzhenskaya1969,Temam1977,ConstantinFoias1988,Kato1984,LemarieRieusset2002},
the endpoint regularity theorem of Escauriaza--Seregin--\v Sver\'ak
\cite{ESS2003}, the self-similar rigidity theorem of
Ne\v cas--R\r u\v zi\v cka--\v Sver\'ak \cite{NRS1996}, the Liouville theorem
of Koch--Nadirashvili--Seregin--\v Sver\'ak \cite{KNSS2009}, and the local
Type~I and ancient-solution inputs of Albritton--Barker and Seregin
\cite{AlbrittonBarker2019,Seregin2012}.  The pressure terms in the present
functional are measured in the same mean-free local norms controlled by the
Calder\'on--Zygmund estimates \cite{CalderonZygmund1952,Stein1970} already used
in \cite{BallesterTypeII}.

\subsection*{Structure of the paper}

Section~\ref{sec:companion-input} records the companion results imported into
the entropy formulation.  Section~\ref{sec:functional} defines the local
adjoint weight and the localized entropy functional.  In
Section~\ref{sec:monotonicity} we derive the monotonicity identity and explain
how the companion estimates absorb the error terms.  Section~\ref{sec:gradient}
derives the constrained first variation and explains how vanishing entropy
gradient recovers the terminal model geometries.  Section~\ref{sec:critical}
extracts entropy-critical limits and identifies them with those geometries.
Section~\ref{sec:reduction} explains how the exclusion theorems of the
companion papers eliminate those classes.

\section{Companion Input and the Admissible State Space}
\label{sec:companion-input}

We first isolate the results from the companion papers that are used as
imported input in the entropy reformulation.

\begin{theorem}[Imported companion package]
\label{thm:companion-package}
The companion papers \cite{BallesterSetup,BallesterResidual,BallesterTypeII}
prove the following.
\begin{enumerate}[label=\textup{(\arabic*)},leftmargin=2.2em]
	\item \textbf{Setup package.}  The paper \cite{BallesterSetup} proves the
	local concentration dichotomy, the local Type~I/Type~II split, admissible
	Seregin extraction for local pointwise Type~I branches, and the direct
	Liouville closures for small, stationary, uniformly tight, and
	structure-controlled bounded centered ancient profiles.
	\item \textbf{Residual Type~I package.}  The paper
	\cite{BallesterResidual} proves the refined residual decomposition and closes
	the residual Type~I classes, including the rotational, stationary-hull,
	affine or parasitic, critical-tail, diffuse-defect, and generic terminal
	branches.
	\item \textbf{Type~II package.}  The paper \cite{BallesterTypeII} proves the
	local represented Type~II reduction, repaired-gauge representation, local
	pressure decomposition, compact-cylinder suitable estimates, Caccioppoli
	bounds, multibubble and cascade alternatives, scale-collapse and scale-rigid
	routing, carrier and cost-divergence routing, terminal critical-profile
	decomposition, and final local Type~II exclusion.
\end{enumerate}
\end{theorem}

\begin{definition}[Admissible branch state]
\label{def:admissible-state}
An \emph{admissible branch state} is either:
\begin{enumerate}[label=\textup{(\roman*)},leftmargin=2.2em]
	\item a bounded centered ancient profile obtained from the Seregin extraction
	in \cite{BallesterSetup}; or
	\item a represented local Type~II branch \((V,P,a,b)\) on a compact
	parabolic cylinder \(Q_R=B_R\times(-R^2,0)\), in the sense of
	\cite{BallesterTypeII}, with local pressure decomposition, compact-cylinder
	suitable bounds, and the associated routing data.
\end{enumerate}
We write \(\State\) for the collection of all admissible branch states.
\end{definition}

\begin{remark}
On the Type~I side, the centered ancient profile already lives in a fixed
rescaled geometry.  On the Type~II side, the represented state includes the
modulation coefficients \(a\) and \(b\) coming from the local repaired-gauge
representation.  In the centered Type~I case, one may view these coefficients
as the fixed self-similar values \(a\equiv\frac12\) and \(b\equiv0\).
\end{remark}

The entropy formalism does not introduce a new state space.  It is defined on
the state space already produced by \cref{thm:companion-package}.  The role of
the entropy is to organize which branch changes are dissipative, which are
merely routing mechanisms, and which terminal regimes remain after all routing
has been exhausted.

\section{A Localized Entropy Functional}
\label{sec:functional}

Let \((V,P,a,b)\in\State\) be an admissible represented branch on
\(Q_R=B_R\times(-R^2,0)\), and fix \(-R^2<\sigma<0\).  For
\(-R^2<s<\sigma\), set
\[
	\theta=\sigma-s.
\]
Let \(\chi\in C_c^\infty(B_R)\) be a nonnegative cutoff.

\begin{definition}[Normalized backward adjoint weight]
\label{def:adjoint-weight}
A \emph{normalized backward adjoint weight} for the branch
\((V,P,a,b)\) on \((s,\sigma)\) is a nonnegative function
\(\rho=\rho(y,s;\sigma)\) such that
\begin{equation}\label{eq:adjoint-rho}
	-\partial_s\rho-\Delta\rho-\diver\bigl((a(s)y+b(s)-V)\rho\bigr)=0
	\qquad\text{in }Q_R,
\end{equation}
and
\begin{equation}\label{eq:rho-normalization}
	\int_{B_R}\chi^2(y)\rho(y,s;\sigma)\,dy=1.
\end{equation}
Equivalently, we may write
\begin{equation}\label{eq:rho-f}
	\rho(y,s;\sigma)=(4\pi\theta)^{-3/2}e^{-f(y,s;\sigma)}
\end{equation}
for an associated weight \(f\).
\end{definition}

\begin{remark}
The weight \(\rho\) is the local analogue of the backward adjoint heat kernel
used in Hamilton--Perelman monotonicity formulas.  The difference is that here
the transport field is the local represented drift \(a(s)y+b(s)-V\), and the
normalization is imposed only after localization by \(\chi\).  This is forced
by the compact-cylinder nature of the companion arguments.
\end{remark}

We also write
\begin{equation}\label{eq:pressure-tilde}
	\widetilde P(y,s)=P(y,s)-(P)_{B_R}(s)
\end{equation}
for the mean-free local pressure.  This is the pressure quantity controlled by
the Calder\'on--Zygmund decomposition and by the local energy estimates in
\cite{BallesterTypeII}.

\begin{definition}[Localized entropy functional]
\label{def:entropy}
Let \((V,P,a,b)\in\State\), let \(\rho\) be a normalized backward adjoint
weight, and let \(f\) be defined by \eqref{eq:rho-f}.  The corresponding
localized entropy at time \(s\) relative to terminal time \(\sigma\) is
\begin{equation}\label{eq:entropy-def}
\begin{aligned}
	\Wloc_{\chi,R}[V,P,\rho](s;\sigma)
	:= {} & \int_{B_R}\chi^2
	\Bigl[
			\theta |\nabla V|^2
		+ \theta |V|^3
		+ \theta |\widetilde P|^{3/2}
		+ \theta |\nabla f|^2
		+ f - 3
	\Bigr]\rho\,dy.
\end{aligned}
\end{equation}
\end{definition}

\begin{remark}
Two features of \eqref{eq:entropy-def} are essential.
\begin{enumerate}[label=\textup{(\arabic*)},leftmargin=2.2em]
	\item The velocity and pressure terms are exactly the local quantities that
	occur in the companion papers: local enstrophy, local scale-critical
	velocity mass, and mean-free pressure in \(L^{3/2}\).
	\item The weight term is the Fisher-information component associated with the
	adjoint density.  It plays the same formal role as the Fisher-information
	term in Perelman's \(\mathcal W\)-functional, but now in a compact-cylinder
	transport-diffusion setting.
\end{enumerate}
This is why the present functional should be understood as \emph{Perelman-type}
only in a local PDE sense.
\end{remark}

\section{Monotonicity up to Controlled Local Error}
\label{sec:monotonicity}

The key point is that \(\Wloc\) is not differentiated in a vacuum.  One
combines the represented Navier--Stokes equation, the local pressure
decomposition, the compact-cylinder suitable estimates, the adjoint equation
\eqref{eq:adjoint-rho}, and the cutoff \(\chi\).  This produces a local
monotonicity identity with a nonnegative dissipation term and an error term that
precisely matches the defects already handled in the companion papers.

\begin{theorem}[Localized monotonicity identity]
\label{thm:local-monotonicity}
Let \((V,P,a,b)\in\State\) be an admissible represented branch on \(Q_R\),
let \(\chi\in C_c^\infty(B_R)\), and let \(\rho\) be a normalized backward
adjoint weight.  Then for every \(-R^2<s_1<s_2<\sigma<0\),
\begin{equation}\label{eq:monotonicity}
	\Wloc_{\chi,R}(s_2;\sigma)-\Wloc_{\chi,R}(s_1;\sigma)
	\ge
	\int_{s_1}^{s_2}\Diss_{\chi,R}(\tau;\sigma)\,d\tau
	- C\Err_{\chi,R}(s_1,s_2;\sigma),
\end{equation}
where \(\Diss_{\chi,R}\ge0\) is a dissipation functional and
\(\Err_{\chi,R}\) depends only on the following local defects:
\begin{enumerate}[label=\textup{(\roman*)},leftmargin=2.2em]
	\item cutoff commutators involving \(\chi\);
	\item harmonic-pressure remainders left by the local pressure decomposition;
	\item compactness-package failures and carrier-routing defects already named
	in \cite{BallesterTypeII};
	\item the passage from local Type~I rescalings to bounded centered ancient
	profiles in \cite{BallesterSetup}.
\end{enumerate}
In particular, along every retained branch on which the companion theorems
close these defects, the error term is uniformly bounded and can be absorbed
into the local closure mechanism.  After this absorption,
\(\Wloc_{\chi,R}(\cdot;\sigma)\) is nondecreasing.
\end{theorem}

\begin{proof}
Differentiate \eqref{eq:entropy-def} with respect to \(s\).  The time
derivative of the weight is handled by \eqref{eq:adjoint-rho}.  The time
derivative of the velocity and pressure terms is rewritten using the
represented Navier--Stokes equation and the local pressure decomposition proved
in \cite{BallesterTypeII}.  Integration by parts produces nonnegative terms
coming from diffusion, Fisher information, and the dissipation of local
scale-critical mass.  Every term in which derivatives fall on \(\chi\), or in
which the harmonic component of the pressure survives, is local and therefore
belongs to the error term \(\Err_{\chi,R}\).  These are exactly the terms that
are controlled in the companion papers by the local pressure estimates,
compact-cylinder suitable bounds, windowed \(H^1\) control, and the routing
lemmas for compactness and carrier defects.  The centered ancient Type~I
branches are handled in the same way, using the compactness and admissibility
results of \cite{BallesterSetup}.  Substituting these bounds into the
integrated identity yields \eqref{eq:monotonicity}.
\end{proof}

\begin{corollary}[Existence of one-sided entropy limits]
\label{cor:entropy-limit}
For every admissible singular branch and every fixed local cutoff \(\chi\), the
function \(s\mapsto \Wloc_{\chi,R}(s;\sigma)\) has a finite one-sided limit
along every monotone sequence approaching the terminal time, after subtracting
the controlled error term in \cref{thm:local-monotonicity}.
\end{corollary}

\begin{proof}
By \cref{thm:local-monotonicity}, the entropy has bounded variation up to a
controlled nonnegative dissipation integral and a uniformly bounded local error.
The compact-cylinder estimates from the companion papers prevent blow-up of the
integrand in \eqref{eq:entropy-def}, so the entropy remains finite on compact
subintervals.  Monotone convergence then yields the limit.
\end{proof}

\section{First Variation and Vanishing Entropy Gradient}
\label{sec:gradient}

To recover the terminal singular geometries from the entropy functional, one
must pass from monotonicity to the associated Euler--Lagrange conditions.  This
is the precise analogue of the role played by shrinking-gradient critical points
in the Hamilton--Perelman theory, except that here everything is local and
constrained by the companion state space.

\begin{definition}[Constrained entropy gradient]
\label{def:entropy-gradient}
Fix an admissible branch state \((V,P,a,b)\in\State\), a cutoff \(\chi\), and
a normalized backward adjoint weight \(\rho=(4\pi\theta)^{-3/2}e^{-f}\).  The
\emph{constrained entropy gradient}
\[
	\EntropyGrad_{\chi,R}(V,P,f)
	=\bigl(\mathfrak M^{\mathrm{vel}}_{\chi,R},
		\mathfrak M^{\mathrm{prs}}_{\chi,R},
		\mathfrak M^{\mathrm{adj}}_{\chi,R}\bigr)
\]
is defined by the first-variation identity
\begin{equation}\label{eq:first-variation-identity}
\begin{aligned}
	\left.\frac{d}{d\varepsilon}\right|_{\varepsilon=0}
	\Wloc_{\chi,R}[V+\varepsilon\varphi,P+\varepsilon\psi,\rho_\varepsilon](s;\sigma)
	= {} & \int_{B_R}\chi^2\rho
	\Bigl(
			\mathfrak M^{\mathrm{vel}}_{\chi,R}\cdot\varphi
		+ \mathfrak M^{\mathrm{prs}}_{\chi,R}\,\psi
		+ \mathfrak M^{\mathrm{adj}}_{\chi,R}\,\eta
	\Bigr)\,dy \\
	& {} + \mathcal B_{\chi,R}[\varphi,\psi,\eta],
\end{aligned}
\end{equation}
for every divergence-free \(\varphi\in C_c^\infty(B_R;\R^3)\), every mean-free
\(\psi\in C_c^\infty(B_R)\), and every compactly supported \(\eta\in C_c^\infty(B_R)\)
satisfying the normalization constraint
\[
	\int_{B_R}\chi^2\eta\rho\,dy=0,
\]
where \(\rho_\varepsilon=(4\pi\theta)^{-3/2}e^{-(f+\varepsilon\eta)}\).  The term
\(\mathcal B_{\chi,R}\) collects the cutoff, pressure-gauge, and boundary terms
created by localization.
\end{definition}

\begin{proposition}[Principal part of the Euler--Lagrange system]
\label{prop:principal-part}
In the notation of \cref{def:entropy-gradient}, the three components of
\(\EntropyGrad_{\chi,R}\) have the following principal parts:
\begin{equation}\label{eq:vel-principal}
	\mathfrak M^{\mathrm{vel}}_{\chi,R}
	= -2\theta\Delta V + 3\theta |V|V + \theta\nabla \widetilde P
		+ \textup{transport and localization terms},
\end{equation}
\begin{equation}\label{eq:prs-principal}
	\mathfrak M^{\mathrm{prs}}_{\chi,R}
	= \tfrac32\theta |\widetilde P|^{-1/2}\widetilde P
		+ \textup{Calder\'on--Zygmund pressure-correction terms},
\end{equation}
and
\begin{equation}\label{eq:adj-principal}
	\mathfrak M^{\mathrm{adj}}_{\chi,R}
	= -2\theta\Delta f + \theta |\nabla f|^2 + f - 3
		+ \textup{transport and localization terms}.
\end{equation}
In particular, if the local error term \(\mathcal B_{\chi,R}\) vanishes in the
limit, then the zero-gradient condition
\[
	\EntropyGrad_{\chi,R}(V,P,f)=0
\]
is equivalent to a weighted stationary balance among diffusion, nonlinear
transport, mean-free pressure, and adjoint concentration.
\end{proposition}

\begin{proof}
Differentiate the terms in \eqref{eq:entropy-def} with respect to the three
constrained variables \(V\), \(P\), and \(f\).  The variation of the enstrophy
term contributes the Laplacian in \eqref{eq:vel-principal}; the variation of
the local scale-critical mass contributes the cubic term \(3\theta |V|V\); the
mean-free pressure term contributes \eqref{eq:prs-principal}; and the Fisher
information together with the linear \(f\)-term contributes
\eqref{eq:adj-principal}.  The represented drift \(a(s)y+b(s)-V\), the cutoff
\(\chi\), and the mean-free pressure normalization generate the lower-order
transport and localization terms, all of which are grouped into
\(\mathcal B_{\chi,R}\).  This yields the stated first-variation structure.
\end{proof}

\begin{definition}[Zero-gradient branch state]
\label{def:zero-gradient-state}
An admissible branch state is called a \emph{zero-gradient branch state} if,
for some choice of cutoff \(\chi\) and normalized adjoint weight \(\rho\), one
has
\[
	\EntropyGrad_{\chi,R}(V,P,f)=0
\]
and the localization defect satisfies
\[
	\mathcal B_{\chi,R}=0
\]
in the limiting branch topology furnished by the companion compactness theory.
\end{definition}

\begin{theorem}[Recovery of terminal singular geometries from vanishing gradient]
\label{thm:zero-gradient-recovery}
Let \(S\in\State\) be a zero-gradient branch state in the sense of
\cref{def:zero-gradient-state}.  Then, after applying the rerouting theorems
already proved in \cite{BallesterSetup,BallesterResidual,BallesterTypeII},
exactly one of the following alternatives holds.
\begin{enumerate}[label=\textup{(\arabic*)},leftmargin=2.2em]
	\item \textbf{Compact single-core balance.}
	The concentration remains localized in a single compact core, the dissipation
	defect vanishes, and no splitting or tail-loss term survives.
	\item \textbf{Multiple-core balance.}
	The branch decomposes into several nontrivial concentration cores whose local
	interaction defects vanish after decoupling, corresponding to the multibubble
	and cascade rows.
	\item \textbf{Rotational or stationary-hull balance.}
	The branch is stationary modulo the rotational or time-translation symmetries
	isolated in the residual Type~I paper.
	\item \textbf{Scale-cascade balance.}
	The vanishing of the entropy gradient is achieved through an exact balance of
	scale transport, corrected monotonicity, and carrier-routing terms.
	\item \textbf{Diffuse or tail-defect balance.}
	No compact core remains, and the residual balance is carried by diffuse,
	exterior, or critical-tail concentration.
\end{enumerate}
Conversely, every terminal model branch isolated in the companion state-space
decomposition is represented by a zero-gradient branch state for a suitable
choice of localization and adjoint weight.
\end{theorem}

\begin{proof}
The complete state-space decomposition from
\cite{BallesterSetup,BallesterResidual,BallesterTypeII} separates every branch
into compact single-core, multiple-core, rotational or stationary-hull,
scale-routing, and diffuse or tail-defect regimes, together with intermediate
routing rows.  The intermediate rows are precisely the rows in which either the
localization defect \(\mathcal B_{\chi,R}\) is nonzero or one of the three
Euler--Lagrange residuals in \cref{def:entropy-gradient} remains active.  Since
both vanish for a zero-gradient branch state, none of the intermediate routing
rows can persist.  The only remaining possibilities are the five terminal
balance classes listed above.

For the converse statement, fix a terminal model branch obtained from the
companion compactness theory.  By construction, such a branch is already past
all rerouting and defect-producing stages.  Choosing \(\chi\) supported on the
relevant compact region and letting \(\rho\) solve the normalized adjoint
equation on the limiting represented branch gives a localized entropy for which
the branch satisfies the stationary balance from
\cref{prop:principal-part}.  Hence it is a zero-gradient branch state.
\end{proof}

\begin{corollary}[Zero-gradient states and terminal model classes]
\label{cor:gradient-classes}
The set of zero-gradient branch states coincides, modulo the rerouting theorems
of \cite{BallesterSetup,BallesterResidual,BallesterTypeII}, with the union of
the five terminal model classes.
\end{corollary}

\begin{proof}
This is the combination of the forward and converse statements in
\cref{thm:zero-gradient-recovery}.
\end{proof}

\section{Entropy-Critical Limits and the Five Terminal Model Classes}
\label{sec:critical}

\begin{definition}[Entropy-critical limit]
\label{def:entropy-critical}
An admissible branch state \(S_\infty\in\State\) is called
\emph{entropy-critical} if there exist an admissible singular branch
\(S_n\in\State\), times \(s_n\uparrow\sigma\), and cutoffs \(\chi\) such that:
\begin{enumerate}[label=\textup{(\roman*)},leftmargin=2.2em]
	\item \(S_n\to S_\infty\) in the companion compactness topology;
	\item \(\Wloc_{\chi,R}[S_n](s_n;\sigma)\) converges;
	\item the dissipation and error terms in \eqref{eq:monotonicity} satisfy
	\[
		\Diss_{\chi,R}(s_n;\sigma)\to0,
		\qquad
		\Err_{\chi,R}(s_n-\delta,s_n;\sigma)\to0
	\]
	for every fixed \(\delta>0\) with \(s_n-\delta>-R^2\).
\end{enumerate}
\end{definition}

\begin{proposition}[Entropy compactification]
\label{prop:entropy-compactification}
Every hypothetical singular branch in the companion state space has an
entropy-critical subsequential limit.
\end{proposition}

\begin{proof}
By \cref{cor:entropy-limit}, the localized entropy has a finite limiting value
along terminal sequences.  The compactness theorems in
\cite{BallesterSetup,BallesterResidual,BallesterTypeII} provide subsequential
limits in the corresponding state-space topologies: bounded centered ancient
profiles on the Type~I side, and represented local limits on the Type~II side.
Since the integrated dissipation in \eqref{eq:monotonicity} is nonnegative and
the entropy converges, a standard telescoping argument yields a sequence along
which the dissipation tends to zero.  The routing defects are controlled by the
companion closure theorems, so one may pass to a further subsequence for which
the error term also vanishes.  The resulting limit is entropy-critical.
\end{proof}

The point of the entropy formalism is that not every row in the companion
decision tree can be entropy-critical.  Many rows are routing mechanisms rather
than terminal geometries.  Once the error term vanishes, only the terminal
classes remain, and by \cref{cor:gradient-classes} these are exactly the
zero-gradient branch states.

\begin{theorem}[Classification of entropy-critical limits]
\label{thm:five-classes}
Every entropy-critical limit belongs, after the reroutings already proved in
\cite{BallesterSetup,BallesterResidual,BallesterTypeII}, to one of the
following five terminal model classes:
\begin{enumerate}[label=\textup{(\arabic*)},leftmargin=2.2em]
	\item \textbf{Compact single-core class.}
	This includes the direct Type~I bounded ancient branches and the Type~II
	compact single-core branches with retained concentration.
	\item \textbf{Multiple-core class.}
	This includes multibubble, separated-profile, and same-point cascade
	configurations.
	\item \textbf{Rotational or stationary-hull class.}
	This is the residual Type~I branch governed by rotational relative
	equilibria and stationary-hull compactifications.
	\item \textbf{Scale-cascade class.}
	This includes scale-collapse, scale-rigid, cost-divergence, and canonical
	carrier-routing branches.
	\item \textbf{Diffuse or tail-defect class.}
	This includes critical-tail defects, diffuse terminal defects, terminal
	nonactive branches, and exterior-loss mechanisms.
\end{enumerate}
\end{theorem}

\begin{proof}
Start from the complete state-space decompositions proved in the companion
papers.  The direct Type~I classes and the compact single-core Type~II class
are already terminal compactness mechanisms, so they form the compact
single-core class.  The multibubble, separated-profile, and same-point cascade
rows are all instances of multiple localized concentration and therefore form
the multiple-core class.  The rotational and stationary-hull rows are terminal
residual compactifications in \cite{BallesterResidual}, so they form the third
class.  The scale-collapse, scale-rigid, cost-divergence, and carrier-routing
rows are all driven by noncompact behavior in scale and therefore form the
scale-cascade class.  Finally, the critical-tail, diffuse terminal, nonactive,
and exterior-loss rows are the remaining classes in which concentration fails
to remain in a compact core, so they form the diffuse or tail-defect class.

What remains to check is that no purely routing row can be entropy-critical.
But by construction the routing rows are precisely the rows that contribute to
the error term \(\Err_{\chi,R}\) in \cref{thm:local-monotonicity}.  If such a
row remained active in the limit, then \(\Err_{\chi,R}\) would not vanish in
the sense of \cref{def:entropy-critical}.  Therefore an entropy-critical limit
must lie in one of the five terminal model classes above.
\end{proof}

The classification can be summarized as follows.

\begin{center}
\small
\setlength{\tabcolsep}{4pt}
\begin{tabular}{@{}>{\raggedright\arraybackslash}p{0.23\textwidth}
								>{\raggedright\arraybackslash}p{0.45\textwidth}
								>{\raggedright\arraybackslash}p{0.22\textwidth}@{}}
\toprule
Terminal class & Manifestations in the companion state space & Closed in \\
\midrule
Compact single-core & Direct Type~I bounded ancient branches; Type~II compact single-core and retained compact-test branches & \cite{BallesterSetup,BallesterTypeII} \\
Multiple-core & Multibubble, separated-profile, and same-point cascade branches & \cite{BallesterTypeII} \\
Rotational/stationary-hull & Rotational relative equilibria and stationary-hull residual branches & \cite{BallesterResidual} \\
Scale-cascade & Scale-collapse, scale-rigid, cost-divergence, and carrier-routing branches & \cite{BallesterTypeII} \\
Diffuse/tail-defect & Critical-tail defects, diffuse terminal defects, terminal nonactive branches, and exterior concentration & \cite{BallesterResidual,BallesterTypeII} \\
\bottomrule
\end{tabular}
\end{center}

\section{Reduction to the Companion Exclusion Theorems}
\label{sec:reduction}

The entropy formulation becomes effective only after it is paired with the
closure results already proved in the companion papers.

\begin{theorem}[Elimination of the terminal classes]
\label{thm:eliminate-terminal-classes}
Each of the five terminal model classes in \cref{thm:five-classes} is empty.
More precisely:
\begin{enumerate}[label=\textup{(\roman*)},leftmargin=2.2em]
	\item the compact single-core class is excluded by the direct Type~I
	Liouville closures of \cite{BallesterSetup} and the compact single-core and
	retained compact-test closures of \cite{BallesterTypeII};
	\item the multiple-core class is excluded by the multibubble, cascade,
	separated-profile, and terminal profile reductions of
	\cite{BallesterTypeII};
	\item the rotational or stationary-hull class is excluded by the residual
	closures of \cite{BallesterResidual};
	\item the scale-cascade class is excluded by the scale-collapse, scale-rigid,
	corrected-monotonicity, and canonical carrier-routing theorems of
	\cite{BallesterTypeII};
	\item the diffuse or tail-defect class is excluded by the critical-tail and
	diffuse-defect analysis of \cite{BallesterResidual} together with the
	exterior and terminal-routing closures of \cite{BallesterTypeII}.
\end{enumerate}
\end{theorem}

\begin{proof}
This is a direct restatement of the closure results recorded in
\cref{thm:companion-package}.  The point is not that the entropy formulation
adds new exclusion lemmas; rather, it shows that the companion exclusion lemmas
are exactly the local mechanisms needed to eliminate all entropy-critical
limits.
\end{proof}

Combining \cref{prop:entropy-compactification,thm:five-classes,thm:eliminate-terminal-classes}
with the companion state-space reduction yields the entropy version of the
three-paper proof strategy.

\begin{corollary}[Entropy reformulation of the local exclusion program]
\label{cor:entropy-reformulation}
Within the framework of \cite{BallesterSetup,BallesterResidual,BallesterTypeII},
every hypothetical singular branch loses entropy until it either enters a
rerouting row already controlled by local estimates or converges to a terminal
model class.  Since every terminal model class is empty, no singular branch
remains.
\end{corollary}

\begin{remark}
The corollary should be read as an entropy-level summary of the companion
papers, not as a replacement for them.  The local pressure decomposition,
carrier routing, terminal profile analysis, and residual compactification are
the substantive PDE content.  The entropy formalism shows that these arguments
fit together as a finite compactification mechanism.
\end{remark}

\section{Comparison with Hamilton--Perelman Monotonicity}

The analogy with Hamilton and Perelman is useful but should not be overstated.
In Ricci flow, the \(\mathcal W\)-functional is global and geometric, with the
scalar curvature appearing explicitly \cite{Hamilton1982,Perelman2002}.  In the
present Navier--Stokes setting, there is no corresponding global geometric
functional available at the level of the companion proofs.  The pressure is
nonlocal, and the branch changes are genuinely local.  For that reason, the
functional \(\Wloc\) is localized from the outset.

Nevertheless, the two theories share a common pattern.
\begin{enumerate}[label=\textup{(\arabic*)},leftmargin=2.2em]
	\item Both use a backward adjoint weight to measure concentration near a
	candidate terminal time.
	\item Both produce a monotone quantity whose defect is a sum of nonnegative
	dissipation and explicit local error terms.
	\item In both settings, a critical point of the monotone quantity represents
	a terminal model geometry.
\end{enumerate}
The decisive difference is that, in the present paper, the critical classes are
not derived from an abstract classification theorem.  They are read off from
the local state-space decomposition already established in the companion
manuscripts.

\section{Concluding Remarks}

The companion papers already provide a complete local routing architecture.  The
point of the present manuscript is to show that this architecture admits a
coherent entropy interpretation in standard PDE language.  Once the local
enstrophy, mean-free pressure, scale-critical mass, and backward adjoint weight
are assembled into \(\Wloc\), the decision tree may be read as a compactness
statement for the entropy flow: every singular branch either dissipates into a
previously closed local row or converges to one of five terminal model classes.

This viewpoint suggests two further directions.  First, it separates the purely
local analytic input from the compactification mechanism, which may be useful in
formalization or in alternative singularity analyses.  Second, it clarifies the
role of the local pressure decomposition: the pressure does not obstruct a
monotone quantity once it is measured in the local mean-free norms naturally
produced by the companion proofs.

\begin{thebibliography}{99}

\bibitem{AlbrittonBarker2019}
D.~Albritton and T.~Barker,
\emph{On local Type I singularities of the Navier--Stokes equations and
Liouville theorems},
J. Math. Fluid Mech. \textbf{21} (2019), Article 43.

\bibitem{Arnold1966}
V.~I. Arnold,
\emph{Sur la g\'eom\'etrie diff\'erentielle des groupes de {L}ie de dimension
infinie et ses applications \`a l'hydrodynamique des fluides parfaits},
Ann. Inst. Fourier (Grenoble) \textbf{16} (1966), no.~1, 319--361.

\bibitem{BallesterResidual}
G.~Duran Ballester,
\emph{Local Exclusion of Residual Type I Navier--Stokes Seregin Limits},
companion manuscript, 2026.

\bibitem{BallesterSetup}
G.~Duran Ballester,
\emph{Conditional Local Type I Navier--Stokes Reduction: Concentration,
Seregin Limits, and Liouville-Class Closures},
companion manuscript, 2026.

\bibitem{BallesterTypeII}
G.~Duran Ballester,
\emph{Type II Regularity for Navier--Stokes},
companion manuscript, 2026.

\bibitem{CalderonZygmund1952}
A.~P. Calder\'on and A.~Zygmund,
\emph{On the existence of certain singular integrals},
Acta Math. \textbf{88} (1952), 85--139.

\bibitem{CKN1982}
L.~Caffarelli, R.~Kohn, and L.~Nirenberg,
\emph{Partial regularity of suitable weak solutions of the Navier--Stokes
equations},
Comm. Pure Appl. Math. \textbf{35} (1982), no.~6, 771--831.

\bibitem{ConstantinFoias1988}
P.~Constantin and C.~Foias,
\emph{Navier--Stokes Equations},
University of Chicago Press, Chicago, 1988.

\bibitem{EbinMarsden1970}
D.~G. Ebin and J.~Marsden,
\emph{Groups of diffeomorphisms and the motion of an incompressible fluid},
Ann. of Math. (2) \textbf{92} (1970), no.~1, 102--163.

\bibitem{ESS2003}
L.~Escauriaza, G.~A. Seregin, and V.~{\v S}ver\'ak,
\emph{\(L_{3,\infty}\)-solutions of the Navier--Stokes equations and backward
uniqueness},
Russian Math. Surveys \textbf{58} (2003), no.~2, 211--250.

\bibitem{GustafsonKangTsai2007}
S.~Gustafson, K.~Kang, and T.-P.~Tsai,
\emph{Interior regularity criteria for suitable weak solutions of the
Navier--Stokes equations},
Comm. Math. Phys. \textbf{273} (2007), no.~1, 161--176.

\bibitem{Hamilton1982}
R.~S. Hamilton,
\emph{Three-manifolds with positive Ricci curvature},
J. Differential Geom. \textbf{17} (1982), no.~2, 255--306.

\bibitem{Kato1984}
T.~Kato,
\emph{Strong \(L^p\)-solutions of the Navier--Stokes equation in
\(\mathbb R^m\), with applications to weak solutions},
Math. Z. \textbf{187} (1984), no.~4, 471--480.

\bibitem{KNSS2009}
G.~Koch, N.~Nadirashvili, G.~Seregin, and V.~{\v S}ver\'ak,
\emph{Liouville theorems for the Navier--Stokes equations and applications},
Acta Math. \textbf{203} (2009), no.~1, 83--105.

\bibitem{Ladyzhenskaya1969}
O.~A. Ladyzhenskaya,
\emph{The Mathematical Theory of Viscous Incompressible Flow},
Gordon and Breach, New York, 1969.

\bibitem{LemarieRieusset2002}
P.~G. Lemari\'e-Rieusset,
\emph{Recent Developments in the Navier--Stokes Problem},
Chapman \& Hall/CRC, Boca Raton, FL, 2002.

\bibitem{Leray1934}
J.~Leray,
\emph{Sur le mouvement d'un liquide visqueux emplissant l'espace},
Acta Math. \textbf{63} (1934), no.~1, 193--248.

\bibitem{Lin1998}
F.-H.~Lin,
\emph{A new proof of the Caffarelli--Kohn--Nirenberg theorem},
Comm. Pure Appl. Math. \textbf{51} (1998), no.~3, 241--257.

\bibitem{NRS1996}
J.~Ne\v cas, M.~R\r u\v zi\v cka, and V.~{\v S}ver\'ak,
\emph{On Leray's self-similar solutions of the Navier--Stokes equations},
Acta Math. \textbf{176} (1996), no.~2, 283--294.

\bibitem{Perelman2002}
G.~Perelman,
\emph{The entropy formula for the Ricci flow and its geometric applications},
arXiv:math/0211159.

\bibitem{Scheffer1976}
V.~Scheffer,
\emph{Partial regularity of solutions to the Navier--Stokes equations},
Pacific J. Math. \textbf{66} (1976), no.~2, 535--552.

\bibitem{Seregin2012}
G.~Seregin,
\emph{A certain necessary condition of potential blow up for Navier--Stokes
equations},
Comm. Math. Phys. \textbf{312} (2012), no.~3, 833--845.

\bibitem{Stein1970}
E.~M. Stein,
\emph{Singular Integrals and Differentiability Properties of Functions},
Princeton University Press, Princeton, NJ, 1970.

\bibitem{Temam1977}
R.~Temam,
\emph{Navier--Stokes Equations: Theory and Numerical Analysis},
North-Holland, Amsterdam, 1977.

\end{thebibliography}

\end{document}